
\chapter{Dualizing Complexes and Local Duality}
\setcounter{section}{-1}

\section{Introduction}\label{sec:5.0}

In this chapter we discuss dualizing complexes on a locally noetherian prescheme $X$. A dualizing complex is a complex $R^\bullet\in\DD^+(X)$ such that the functor
\[
\mathscr{D}\colon M^{\bullet}\longmapsto \RsheafHom^\bullet(M^{\bullet},R^{\bullet})
\]
induces an auto-duality of the category $\DD_{\mathrm{c}}^{b}(X)$ consisting of those bounded complexes in $\DD^+(X)$ which have coherent cohomology. We will show that a large class of preschemes admits dualizing complexes, and that they are almost uniquely determined.

The notion of dualizing complex will allow us to write the duality theorem in a new way. For example, let $X$ be projective $n$-space over a field $k$. Then $R^\bullet=\omega[n]$ (where $\omega=\omega_{X/k}$ is the sheaf of relative $n$-differentials \bookref{sec:3.1}{III,~\S1}) is a dualizing complex for $X$, and $k$ is a dualizing complex for $k$. We have a canonical isomorphism $\RR f_*(\omega[n])\cong k$ \bookref{thm:3.3.4}{III,~3.4}, and hence, for any complex $F^\bullet\in\DD^b_{\mathrm{c}}(X)$, a homomorphism
\[
\Theta\colon\ \RR f_* \mathscr{D}_X(F^\bullet)\longrightarrow \mathscr{D}_k(\RR f_*(F^\bullet)),
\]
where $\mathscr{D}_X$ and $\mathscr{D}_k$ are the dualizing functors on $X$ and on $k$, respectively. The duality theorem \bookref{thm:3.5.1}{III,~5.1} says that $\Theta$ is an isomorphism, i.e., that the dualizing functors on $X$ and on $k$ commute with $\RR f_*$.

In the latter part of the chapter we discuss the local duality theorem, as in \cite[\nopp \S\S2,4,6]{LC} or \cite[\nopp 62, expos\'es~2,4]{SGA}. The proof of the duality theorem given here is different, however.

The reader will see that the existence of a dualizing complex on a locally noetherian prescheme $X$ implies that $X$ has finite Krull dimension. Hence we will give a more general notion, that of a pointwise dualizing complex, to cover the case of infinite Krull dimension.

For convenience we defer the question of existence of dualizing complexes until the end of the chapter (section~\hyperref[sec:5.10]{10}).

\section{Example: duality for abelian groups}\label{sec:5.1}

Let $X=\operatorname{Spec}\mathbb{Z}$. Then the category of quasi-coherent sheaves on $X$ is isomorphic to the category $(\Ab)$ of abelian groups. We know that $\mathbb{Q}/\mathbb{Z}$ gives a good duality for finite abelian groups, i.e., that the functor $\operatorname{Hom}(\cdot,\mathbb{Q}/\mathbb{Z})$ is an exact contravariant functor, which, when applied twice to a finite abelian group, gives that group back again. Similarly, the group $\mathbb{Q}$ gives a good notion of duality for finitely generated free abelian groups. Combining the two, we consider the complex
\[
\mathbb{Q}\longrightarrow \mathbb{Q}/\mathbb{Z}
\]
and work in the derived category $\DD_{\mathrm{c}}^{b}(\Ab)$, consisting of those bounded complexes in $\DD^{+}(\Ab)$ which have finitely generated cohomology. ($\mathrm{c}$ is for ``coherent'', which means ``finitely generated'' in our case.) Since the complex above is an injective resolution of $\mathbb{Z}$, hence isomorphic to $\mathbb{Z}$ in the derived category, we arrive at the following proposition.

\begin{bookproposition}{1.1}\label{prop:5.1.1}
The functor
\[
D\colon M^{\bullet}\longmapsto \RHom^\bullet(M^{\bullet},\mathbb{Z})
\]
is a contravariant $\partial$-functor from $\DD_{\mathrm{c}}^{b}(\Ab)$ into itself, and there is a natural functorial isomorphism
\[
\eta\colon \id_{\DD_{\mathrm{c}}^{b}(\Ab)}\xrightarrow{\sim} DD.
\]
\end{bookproposition}

\begin{proof}
Since the cohomology of $D(M^{\bullet})$ is $\operatorname{Ext}^{i}(M^{\bullet},\mathbb{Z})$, if $M^{\bullet}\in\DD^{b}_{\mathrm{c}}(\Ab)$, then so is $D(M^{\bullet})$. The natural functorial homomorphism $\eta$ is defined in Lemma~\ref{lem:5.1.2} below. To show that it is an isomorphism, by taking a free resolution of $M^{\bullet}$ and applying the lemma on way-out functors \bookref{prop:1.7.1}{I,~7.1}, we reduce to the case $M^{\bullet}=\mathbb{Z}^{r}$ for some $r$. Since the functors in question are additive, one reduces to the case $M^{\bullet}=\mathbb{Z}$, and to complete the proof, one need only observe that $\operatorname{Ext}^{i}(\mathbb{Z},\mathbb{Z})=\mathbb{Z}$ if $i=0$, and $0$ otherwise.
\end{proof}

\begin{booklemma}{1.2}\label{lem:5.1.2}
Let $X$ be a prescheme, let $R^\bullet\in\DD^+(X)$ and let $F^\bullet\in\DD(X)$. Then there is a natural functorial homomorphism
\[
\eta\colon F^\bullet\longrightarrow \RsheafHom^\bullet(\RsheafHom^\bullet(F^\bullet,R^\bullet),R^\bullet).
\]
\end{booklemma}

\begin{proof}
Replacing $R^\bullet$ by an injective resolution, we can erase the $\RR$. As usual, we need only define $\eta$ on objects of $\DD(X)$, i.e., complexes. Given an index $p$, and a section $s\in F^{p}(U)$ over an open set $U$ in $X$, we must define, for each $q$, a homomorphism of sheaves
\[
\sheafHom^{q}(F^\bullet,R^\bullet)_U\longrightarrow R_U^{p+q},
\]
where the subscript $U$ denotes restriction of sheaves to $U$. Given a homomorphism $f\in\sheafHom^{q}(F^\bullet,R^\bullet)_U(V)$ defined on an open set $V\subseteq U$, we send it to the section $f(s_V)\in R^{p+q}(V)$, where $s_V$ is the section $s$, restricted to $V$. One checks that this indeed does define a morphism of complexes, and that it is functorial in $F^\bullet$.
\end{proof}

\section{Dualizing complexes}\label{sec:5.2}
\begingroup
\setlength{\emergencystretch}{2em}

Throughout this section, $X$ will denote a locally noetherian prescheme. We will consider complexes $R^\bullet\in\DD^+(X)$ which have coherent cohomology and finite injective dimension, i.e., $R^\bullet\in\DD_{\mathrm{c}}^+(X)_{\mathrm{fid}}$ (see \bookref{prop:1.7.6}{I,~7.6} and \bookref{prop:2.7.20}{II,~7.20}). Note that if $X$ is quasi-compact, then the condition of finite injective dimension on a complex $R^\bullet\in\DD_{\mathrm{c}}^+(X)$ can be checked locally (e.g., by \bookref{prop:2.7.20}{II,~7.20(iii)}). Note also that such a complex $R^\bullet$ is quasi-isomorphic to a bounded complex of quasi-coherent injective sheaves on $X$ \bookref{prop:2.7.20}{II,~7.20(i)$_{\mathrm{qc}}$}, and hence (using \bookref{prop:2.3.3}{II,~3.3}) the functor
\[
\mathscr{D}\colon F^{\bullet}\longmapsto \RsheafHom^{\bullet}(F^{\bullet}, R^{\bullet})
\]
sends $\DD_{\mathrm{c}}^{b}(X)$ into itself, and interchanges $\DD_{\mathrm{c}}^{+}(X)$ and $\DD_{\mathrm{c}}^{-}(X)$.

\begin{definition*}
Let $R^\bullet\in\DD^+(X)$. We say that a complex $F^\bullet\in\DD(X)$ is \emph{reflexive} (with respect to $R^\bullet$) if the natural map
\[
\eta\colon F^{\bullet}\longrightarrow \RsheafHom^{\bullet}(\RsheafHom^{\bullet}(F^{\bullet}, R^{\bullet}), R^{\bullet})
\]
of Lemma~\ref{lem:5.1.2} is an isomorphism.
\end{definition*}

\begin{bookproposition}{2.1}\label{prop:5.2.1}
Let $R^\bullet\in\DD^+_{\mathrm{c}}(X)_{\mathrm{fid}}$. Then the following conditions are equivalent:
\begin{itemize}
\item[(i)] Every $F^\bullet\in\DD_{\mathrm{c}}(X)$ is reflexive (with respect to $R^\bullet$).
\item[(ii)] Every $F^\bullet\in\DD_{\mathrm{c}}^b(X)$ is reflexive.
\item[(iii)] Every coherent sheaf on $X$ is reflexive.
\item[(iv)] The structure sheaf $\sO_X$ is reflexive.
\end{itemize}
\end{bookproposition}

\begin{proof}
Clearly (i) $\Longrightarrow$ (ii) $\Longrightarrow$ (iii) $\Longrightarrow$ (iv). To prove (iv) $\Longrightarrow$ (i) we note that the question is local on $X$. Hence we may assume that $X$ is affine. Then every coherent sheaf on $X$ is the quotient of a free sheaf $\sO_X^r$; the functors in question are additive, so the result follows from the Lemma on Way-Out Functors \bookref{prop:1.7.1}{I,~7.1(ii) and (iv)}, taking into account \bookref{prop:2.7.20}{II,~7.20(ii)} which implies that $\RsheafHom^\bullet(\cdot,R^\bullet)$ is way-out in both directions.
\end{proof}

\begin{definition*}
Let $X$ be a locally noetherian prescheme. A complex $R^\bullet\in\DD^+_{\mathrm{c}}(X)_{\mathrm{fid}}$ satisfying the equivalent conditions of the proposition is called a \emph{dualizing complex} for $X$.
\end{definition*}

\begin{definition*}
The \emph{Krull dimension} of a locally noetherian prescheme $X$ is the largest $n$ (or $+\infty$) for which there is a chain $Z_0<Z_1<\cdots<Z_n$ of closed irreducible subsets $Z_i$ of $X$. It is also the sup of the dimensions of the local rings of $X$.
\end{definition*}

\begin{bookexample}{2.2}\label{ex:5.2.2}
Let $X$ be a regular prescheme of finite Krull dimension. (A \emph{regular prescheme} is a prescheme $X$ such that all the local rings $\sO_{X,x}$ of points of $X$ are regular local rings.) Then the structure sheaf $\sO_X$ is a dualizing complex for $X$.

Indeed, condition (iv) of the proposition is trivial, because $\sheafExt^{i}(\sO_X,\sO_X)=\sO_X$ for $i=0$, and $0$ otherwise. Clearly $\sO_X\in\DD_{\mathrm{c}}^{+}(X)$. To show that $\sO_X$ has finite injective dimension, we use condition $(\underline{\mathrm{iii}})_{\mathrm{c}}$ of \bookref{prop:2.7.20}{II,~7.20}. Let $n_0$ be the Krull dimension of $X$. Then for every coherent sheaf $F$ on $X$, $\sheafExt^{i}(F,\sO_X)=0$ for $i>n_0$, because $\sheafExt^{i}$ commutes with taking stalks, and the local rings $\sO_{X,x}$ all have cohomological dimension $\leq n_0$, since they are regular local rings \cite[\nopp (28.2)]{14}.
\end{bookexample}

\begin{bookcorollary}{2.3}\label{cor:5.2.3}
Let $R^\bullet\in\DD_{\mathrm{c}}^+(X)_{\mathrm{fid}}$. Then $R^\bullet$ is dualizing if and only if for every $x\in X$, $R_x^\bullet$ is dualizing on the local scheme $\operatorname{Spec}\sO_{X,x}$. Furthermore, it is sufficient to take the closed points $x\in X$.
\end{bookcorollary}

\begin{proof}
For $x\in X$, we see that the stalk $R_x^\bullet$ has finite injective dimension in $\DD^+_{\mathrm{c}}(\operatorname{Spec}\sO_{X,x})$ by using \bookref{prop:2.7.20}{II,~7.20(i)$_{\mathrm{qc}}$} and \bookref{prop:2.7.17}{II,~7.17(iii)}. Thus the result follows from condition (iv) of the Proposition.
\end{proof}

\begin{bookproposition}{2.4}\label{prop:5.2.4}
Let $f\colon X\to Y$ be a finite morphism of locally noetherian preschemes, and let $R^\bullet$ be a dualizing complex on $Y$. Then $f^{\flat}R^\bullet$ is a dualizing complex on $X$ (cf.\ \bookref{sec:3.6}{III,~\S6}).
\end{bookproposition}

\begin{proof}
We may assume that $R^\bullet$ is a bounded complex of quasi-coherent injective sheaves on $Y$. Then
\[
f^{\flat}(R^{\bullet}) = \sheafHom_{Y}(f_{*}\sO_{X},R^{\bullet})^{\sim},
\]
which is a bounded complex of quasi-coherent injectives on $X$. Hence $f^{\flat}(R^{\bullet})\in\DD^{b}_{\mathrm{c}}(X)_{\mathrm{fid}}$. To see that $f^{\flat}(R^{\bullet})$ is dualizing, we apply condition (iv) of Proposition~\ref{prop:5.2.1}, and consider the map
\[
\eta\colon \sO_{X}\longrightarrow \RsheafHom_{X}^{\bullet}(\RsheafHom_{X}^{\bullet}(\sO_{X}, f^{\flat}(R^{\bullet})), f^{\flat}(R^{\bullet})).
\]
Applying $\RR f_*$ to this map, and using duality for a finite morphism \bookref{thm:3.6.7}{III,~6.7} twice, we see that $\RR f_*(\eta)$ is an isomorphism. But we are dealing with quasi-coherent sheaves, and $f_*$ is faithful, so $\eta$ is an isomorphism.
\end{proof}

\par\medskip\noindent\textbf{Remarks.}
\begin{enumerate}
\item We will see later that if $f$ is a smooth morphism, then $f^{\sharp}$ of a dualizing complex is dualizing (Theorem~\hyperref[thm:5.8.3]{8.3} below).
\item This result, joined with the example above, shows that any closed subscheme of a regular prescheme of finite Krull dimension admits a dualizing complex.
\end{enumerate}
\par\medskip

\begin{bookproposition}{2.5}\label{prop:5.2.5}
A complex $R^\bullet\in\DD_{\mathrm{c}}^+(X)_{\mathrm{fid}}$ is dualizing if and only if for every closed point $x\in X$, the sheaf $k(x)$, consisting of the residue field $k(x)$ at the point $x$ and zero elsewhere, is reflexive.
\end{bookproposition}

\begin{proof}
1) Using Corollary~\ref{cor:5.2.3}, we reduce to the case where $X$ is the spectrum of a local noetherian ring $A$, with residue field $k$. The assumption is then that $k$ is reflexive, and we wish to show for every $A$-module of finite type $M$ (in particular, for $M=A$), that $M$ is reflexive on $X$.

2) Using induction on the length, we show first that every $A$-module $M$ of finite length is reflexive. Indeed, for $\operatorname{length}M=1$, we have assumed it. For $\operatorname{length}M>1$, one can write a short exact sequence
\[
0\longrightarrow M'\longrightarrow M\longrightarrow M''\longrightarrow 0
\]
where $M'$ and $M''$ have length $<\operatorname{length}M$. Then by the induction hypothesis $M'$ and $M''$ are reflexive, and by long exact sequences ($\mathscr{D}\mathscr{D}$ is a $\partial$-functor!) and the five-lemma, $M$ is reflexive.

3) Now we show that any module of finite type $M$ is reflexive, using induction on the dimension of $M$. If $\dim M=0$, then $M$ has finite length, and $M$ is reflexive by 2) above. Let $M$ be a module of finite type. Let $M'$ be the submodule of elements with support $\mathfrak{m}$, where $\mathfrak{m}$ is the maximal ideal of $A$. Then $M'$ has finite length, so it is sufficient to consider $M/M'$. In other words, we have reduced to the case $\operatorname{Ass}M\cap\{\mathfrak{m}\}=\emptyset$.

Let $t\in\mathfrak{m}$ be an element which is not a zero-divisor in $M$, so that we have an exact sequence
\[
0\longrightarrow M\xrightarrow{t} M\longrightarrow M/tM\longrightarrow 0.
\]
Now $\dim M/tM<\dim M$, so by the induction hypothesis $M/tM$ is reflexive. Therefore, for each $i\neq0$, we have an exact sequence
\[
H^{i}(\mathscr{D}\mathscr{D}(M))\xrightarrow{t}H^{i}(\mathscr{D}\mathscr{D}(M))\longrightarrow 0.
\]
Now these $H^{i}$ are $A$-modules of finite type, and $t\in\mathfrak{m}$, so by Nakayama's lemma, $H^{i}(\mathscr{D}\mathscr{D}(M))=0$ for $i\neq0$.

Consider the commutative diagram
\[
\begin{tikzcd}
0\ar[r] & M \ar[r,"t"] \ar[d] & M \ar[r] \ar[d,"\alpha"] & M/tM \ar[r] \ar[d,"\cong"] & 0\\
0\ar[r] & H^0(\mathscr{D}\mathscr{D}(M)) \ar[r,"t"] & H^0(\mathscr{D}\mathscr{D}(M)) \ar[r] & H^0(\mathscr{D}\mathscr{D}(M/tM)) \ar[r] & 0
\end{tikzcd}
\]
with exact rows. A diagram chase shows that
\[
H^{0}(\mathscr{D}\mathscr{D}(M))=\alpha(M)+tH^{0}(\mathscr{D}\mathscr{D}(M)),
\]
so by Nakayama's Lemma, $\alpha$ is surjective. To show that $\alpha$ is injective, let $0\neq x\in M$. Choose $n$ so large that $x\notin t^nM$, and draw the same diagram as above, but with $t^n$ in place of $t$. Then $x$ does not become $0$ in $M/t^nM$, and $M/t^nM\cong H^0(\mathscr{D}\mathscr{D}(M/t^nM))$, so $\alpha(x)\neq0$.

Thus $M$ is reflexive. Taking $M=A$, we find that $R^\bullet$ is a dualizing complex.
\end{proof}

\begin{bookproposition}{2.6}\label{prop:5.2.6}
Let $X$ be a locally noetherian prescheme, let $R^\bullet\in\DD^+_{\mathrm{c}}(X)_{\mathrm{fid}}$ be a dualizing complex, and let $\mathscr{D}$ be the functor $\RsheafHom^\bullet(\cdot,R^\bullet)$. Then
\begin{enumerate}
\item[a)] A complex $G^\bullet\in\DD_{\mathrm{c}}^{\mathrm{b}}(X)$ has finite Tor-dimension \bookref{prop:2.4.2}{II,~4.2} if and only if $\mathscr{D}(G^\bullet)$ has finite injective dimension.
\item[b)] There is a functorial isomorphism
\[
\mathscr{D}(\RsheafHom^\bullet(F^\bullet,G^\bullet)) \xrightarrow{\sim} F^\bullet\derivedtensor \mathscr{D}(G^\bullet)
\]
for $F^\bullet\in\DD_{\mathrm{c}}^-(X)$ and $G^\bullet\in\DD^+(X)$, or for $F^\bullet\in\DD_{\mathrm{c}}(X)$ and $G^\bullet\in\DD^+(X)_{\mathrm{fid}}$.
\item[c)] There is a functorial isomorphism
\[
\mathscr{D}(F^{\bullet}\derivedtensor G^{\bullet}) \xrightarrow{\sim} \RsheafHom^{\bullet}(F^{\bullet}, \mathscr{D}(G^{\bullet}))
\]
for $F^\bullet$ and $G^\bullet\in\DD_{\mathrm{c}}^-(X)$, or for $F^\bullet\in\DD_{\mathrm{c}}(X)$ and $G^\bullet\in\DD_{\mathrm{c}}^{\mathrm{b}}(X)_{\mathrm{fTd}}$.
\end{enumerate}
\end{bookproposition}

\begin{proof}
The natural map of sheaves
\[
F\otimes\sheafHom(G,R)\longrightarrow \sheafHom(\sheafHom(F,G),R)
\]
gives rise to a morphism of functors on the derived category,
\[
\tag{*}\label{eq:5.2.comparison}
F^\bullet\derivedtensor \RsheafHom^\bullet(G^\bullet,R^\bullet)\longrightarrow \RsheafHom^\bullet(\RsheafHom^\bullet(F^\bullet,G^\bullet),R^\bullet),
\]
provided either
\begin{itemize}
\item[1)] $F^\bullet\in\DD^{-}(X)$, $G^\bullet\in\DD^{+}(X)$ and $R^\bullet\in\DD^{+}(X)_{\mathrm{fid}}$, or
\item[2)] $F^\bullet\in\DD(X)$, $G^\bullet,R^\bullet\in\DD^+(X)$, and $\mathscr{D}(G^\bullet)\in\DD^b(X)_{\mathrm{fTd}}$.
\end{itemize}
If under the set of conditions 1), we assume furthermore that $F^\bullet\in\DD^-_{\mathrm{c}}(X)$, then by \bookref{prop:1.7.1}{I,~7.1}, the morphism \eqref{eq:5.2.comparison} is an isomorphism. Indeed, we reduce to the case $F=\sO_X$, which is trivial. Taking the inverse isomorphism gives b) under the first set of hypotheses. As a corollary, we see that if $G^\bullet\in\DD^+(X)$ has finite injective dimension, then $\mathscr{D}(G^\bullet)$ has finite Tor-dimension. Indeed, by the Remark following \bookref{prop:2.4.2}{II,~4.2} it is enough to show that there is an integer $n_0$ such that $\sheafTor_i(F,\mathscr{D}(G^\bullet))=0$ for all $i>n_0$ and all coherent sheaves $F$ on $X$, and this follows from b). Thus we have established one half of a). Now if $G^\bullet\in\DD^+(X)_{\mathrm{fid}}$, then as just remarked, $\mathscr{D}(G^\bullet)\in\DD^b(X)_{\mathrm{fTd}}$, so by 2) and \bookref{prop:1.7.1}{I,~7.1} again, we obtain the isomorphism b) under the second set of conditions.

For statement c), let $F^\bullet,G^\bullet\in\DD^-_{\mathrm{c}}(X)$, apply b) to $F^\bullet$ and $\mathscr{D}(G^\bullet)$, and apply $\mathscr{D}$ to the resulting isomorphism. Then since $\mathscr{D}^2=\id$ on $\DD^-_{\mathrm{c}}(X)$, we have the required isomorphism under the first set of conditions. As a corollary we deduce (using \bookref{prop:2.7.20}{II,~7.20 $(\underline{\mathrm{iii}})_{\mathrm{c}}$}) that if $G^\bullet\in\DD^{\mathrm{b}}_{\mathrm{c}}(X)$ has finite Tor-dimension, then $\mathscr{D}(G^\bullet)$ has finite injective dimension. This gives the other half of a), and applying $\mathscr{D}$ to the second half of b) gives the second half of c).
\end{proof}

\par\endgroup

\section{Uniqueness of the dualizing complex}\label{sec:5.3}
\begingroup
\setlength{\emergencystretch}{2em}

\begin{booktheorem}{3.1}\label{thm:5.3.1}
Let $X$ be a connected, locally noetherian prescheme, and let $R^\bullet$ be a dualizing complex on $X$. Let $R^{\prime\bullet}\in\DD_{\mathrm{c}}^+(X)$ be any other complex. Then $R^{\prime\bullet}$ is a dualizing complex if and only if there exists an invertible sheaf $L$ on $X$, and an integer $n$, such that
\[
R^{\prime\bullet}\cong R^\bullet\derivedtensor L[n].
\]
Furthermore, $L$ and $n$ are determined uniquely to within isomorphism.
\end{booktheorem}

\begin{proof}
1) Observe that such an $R^{\prime\bullet}$ is a dualizing complex. Indeed, we may assume that $R^\bullet$ is a bounded complex of quasi-coherent injectives. Then $R^\bullet\otimes L[n]$ is also a bounded complex of injectives, since the property of being injective is local \bookref{lem:2.7.16}{II,~7.16}, and $L$ is locally free. $R^{\prime\bullet}$ is dualizing by Corollary~\ref{cor:5.2.3}.

2) Observe that $L$ and $n$ are uniquely determined by the identity $L[n]\cong \RsheafHom^\bullet(R^\bullet,R^{\prime\bullet})$.

3) Now we show the existence of $L$ and $n$, assuming that $R^{\prime\bullet}$ is a dualizing complex. Let $\mathscr{D}$ and $\mathscr{D}'$ be the dualizing functors corresponding to $R^\bullet$ and $R^{\prime\bullet}$. Define
\[
L^\bullet = \mathscr{D}'\mathscr{D}(\sO_X) = \RsheafHom^\bullet(R^\bullet,R^{\prime\bullet}).
\]
We will prove that $L^\bullet$ is isomorphic, in $\DD^+(X)$, to a complex of the form $L[n]$, where $L$ is an invertible sheaf, and that $R^{\prime\bullet}\cong R^{\bullet}\derivedtensor L^{\bullet}$.

\begin{booklemma}{3.2}\label{lem:5.3.2}
There is a natural functorial isomorphism, for all $M^{\bullet}\in\DD_{\mathrm{c}}^{-}(X)$,
\[
M^{\bullet}\derivedtensor L^{\bullet}\xrightarrow{\sim} \mathscr{D}'\mathscr{D}(M^{\bullet}).
\]
(Note that this tensor product makes sense since $L^\bullet\in\DD^-_{\mathrm{c}}(X)$, cf. \bookref{sec:2.4}{II,~\S4}.)
\end{booklemma}

\begin{proof}
As in Lemma~\ref{lem:5.1.2} above, one defines a functorial homomorphism
\[
M^{\bullet}\derivedtensor \RsheafHom^{\bullet}(R^{\bullet},R^{\prime\bullet})\xrightarrow{} \RsheafHom^{\bullet}(\RsheafHom^{\bullet}(M^{\bullet},R^{\bullet}),R^{\prime\bullet}).
\]
Since it is an isomorphism for $M^{\bullet}=\sO_X$, by the lemma on way-out functors \bookref{prop:1.7.1}{I,~7.1}, it is an isomorphism for all $M^{\bullet}\in\DD_{\mathrm{c}}^{-}(X)$.
\end{proof}

We now define $L^\bullet = \mathscr{D}'\mathscr{D}(\sO_X)$, and use the lemma to deduce that
\[
L^\bullet \derivedtensor L^{\prime\bullet} = \mathscr{D}'\mathscr{D}(L^{\prime\bullet}) = \mathscr{D}'\mathscr{D}\mathscr{D}\mathscr{D}'(\sO_X) = \sO_X
\]
where $L^{\prime\bullet}=\mathscr{D}\mathscr{D}'(\sO_X)$.

\begin{booklemma}{3.3}\label{lem:5.3.3}
Let $X$ be a connected locally noetherian prescheme, and let $L^\bullet,L^{\prime\bullet}\in\DD_{\mathrm{c}}^{-}(X)$ be such that $L^\bullet\derivedtensor L^{\prime\bullet}\cong\sO_X$. Then $L^\bullet=L[n]$ for some invertible sheaf $L$, and $L^{\prime\bullet}=L^{\vee}[-n]$.
\end{booklemma}

\begin{proof}
Since $X$ is connected, we reduce immediately to the case where $X$ is the spectrum of a local noetherian ring $A$, and we wish to show that $L^\bullet=\sO_X[n]$ for some $n$. Let $p$ be the largest integer such that $\mathcal{H}^p(L^\bullet)\neq0$, and let $q$ be the largest integer such that $\mathcal{H}^q(L^{\prime\bullet})\neq0$. Then one sees easily that $p+q$ is the largest integer for which $\mathcal{H}^n(L^\bullet\derivedtensor L^{\prime\bullet})\neq0$, and $\mathcal{H}^{p+q}(L^\bullet\derivedtensor L^{\prime\bullet})=\mathcal{H}^p(L^\bullet)\otimes \mathcal{H}^q(L^{\prime\bullet})$. Therefore $p+q=0$, and $\mathcal{H}^p(L^\bullet)\otimes \mathcal{H}^q(L^{\prime\bullet})=\sO_X$. It follows, by a simple lemma on modules of finite type over a local noetherian ring (cf.\ \cite[\nopp 0,~5.4.3]{EGA}), that $\mathcal{H}^p(L^\bullet)\cong \mathcal{H}^q(L^{\prime\bullet})\cong\sO_X$. Now since $\sO_X$ is locally free, one shows that $L^\bullet = L_1^\bullet\oplus\sO_X[-p]$, where $L_1^\bullet$ has cohomology only in dimensions $<p$. Similarly, $L^{\prime\bullet}=L_1^{\prime\bullet}\oplus\sO_X[p]$. Thus from $L^\bullet\derivedtensor L^{\prime\bullet}\cong\sO_X$ we deduce
\[
L_{1}^{\bullet}[p]\oplus L_{1}^{\prime\bullet}[-p]\oplus (L_{1}^{\bullet}\derivedtensor L_{1}^{\prime\bullet}) = 0
\]
from which it follows that $L_1^\bullet$ and $L_1^{\prime\bullet}$ are $0$ in $\DD_{\mathrm{c}}(X)$, so that $L^\bullet\cong\sO_X[-p]$, and $L^{\prime\bullet}\cong\sO_X[p]$, as required.
\end{proof}

This proves the lemma, so we have only to show that $R^{\prime\bullet} = R^\bullet\derivedtensor L^{\bullet}$. Indeed, $R^\bullet\in\DD^b_{\mathrm{c}}(X)$, so applying Lemma~\ref{lem:5.3.2},
\[
R^{\bullet}\derivedtensor L^{\bullet} = \mathscr{D}'\mathscr{D}(R^{\bullet}) = \mathscr{D}'(\sO_X) = R^{\prime\bullet}
\]
\end{proof}

We now give an application of the uniqueness theorem.

\begin{bookproposition}{3.4}\label{prop:5.3.4}
Let $A$ be a noetherian local ring with residue field $k$, and let $R^\bullet\in\DD^+_{\mathrm{c}}(A)$. Then $R^\bullet$ is dualizing if and only if there is an integer $d$ such that
\[
\operatorname{Ext}^{i}(k,R^{\bullet}) = \begin{cases} 0 & \text{for } i\neq d \\ k & \text{for } i = d. \end{cases}
\]
(Here we write $\DD^+(A)$ for the derived category of the category of $A$-modules. The subscript $\mathrm{c}$ denotes complexes whose cohomology modules are of finite type. We carry over the definitions and results of the previous sections to this case. (Cf.\ \bookref{cor:2.7.19}{II,~7.19} which ensures that we will not get into trouble.))
\end{bookproposition}

\begin{proof}
First suppose that $R^\bullet$ is dualizing, and let $j\colon\operatorname{Spec}k\longrightarrow\operatorname{Spec}A$ be the inclusion. Then by Proposition~\ref{prop:5.2.4},
\[
j^{\flat}(R^{\bullet}) = \RHom_A^{\bullet}(k,R^{\bullet})
\]
is a dualizing complex on $k$. But $k$ itself is a dualizing complex on $k$, so by the Uniqueness Theorem~\ref{thm:5.3.1} above, $j^{\flat}(R^{\bullet})\cong k[-d]$ for some $d$. This $d$ will do.

For the converse, it is clear that $k$ is reflexive, so by Proposition~\ref{prop:5.2.5} we have only to show that $R^\bullet$ has finite injective dimension. Indeed, we will show by induction on $r=\dim M$, for any $A$-module $M$ of finite type, that
\[
\operatorname{Ext}^{i}(M,R^{\bullet}) = 0 \qquad\text{for } i\notin[d-r,d].
\]
Indeed, for $\dim M=0$, $M$ is of finite length, and our statement follows from the case $M=k$ by induction on the length of $M$.

For $\dim M=r>0$, we induct as in the proof of Proposition~\ref{prop:5.2.5}. First we may assume that $\mathfrak{m}\notin\operatorname{Ass}M$, where $\mathfrak{m}$ is the maximal ideal of $A$. Then choosing $t\in\mathfrak{m}$ with $t$ a non-zero-divisor in $M$, we have
\[
0\longrightarrow M\xrightarrow{t} M\longrightarrow M/tM\longrightarrow 0
\]
and $\dim M/tM<r$. Then for $i\notin[d-r+1,d]$ we have $\operatorname{Ext}^i(M/tM,R^\bullet)=0$, hence for $i\notin[d-r,d]$ we have from the long exact sequence of $\operatorname{Ext}$'s,
\[
\operatorname{Ext}^{i}(M,R^{\bullet})\xrightarrow{t}\operatorname{Ext}^{i}(M,R^{\bullet})\longrightarrow 0.
\]
It follows by Nakayama's lemma that $\operatorname{Ext}^{i}(M,R^{\bullet})=0$.
\end{proof}

\begin{bookcorollary}{3.5}\label{cor:5.3.5}
Let $A$ be a noetherian local ring, and let $R^\bullet\in\DD^+_{\mathrm{c}}(A)$. Then $R^\bullet$ is dualizing on $A$ if and only if $R^\bullet\otimes_{A}\widehat{A}$ is dualizing on the completion $\widehat{A}$ of $A$.
\end{bookcorollary}

\begin{proof}
Clearly $R^\bullet\otimes_{A}\widehat{A}\in\DD_{\mathrm{c}}^{+}(\widehat{A})$. Furthermore,
\[
\operatorname{Ext}_{\widehat{A}}^{i}(k, R^{\bullet}\otimes_{A}\widehat{A}) = \operatorname{Ext}_{A}^{i}(k, R^{\bullet})\otimes_{A}\widehat{A}
\]
for all $i$. Thus the ``only if'' implication is clear. For the ``if'' implication, note that the functor $\otimes_{A}\widehat{A}$ is faithfully flat, so if $M$ is a non-zero $A$-module, then $M\otimes_{A}\widehat{A}$ is non-zero. Furthermore, if $M$ is an $A$-module, such that $M\otimes_{A}\widehat{A}\cong k$, then $M\cong k$. (Indeed, map $M$ to $k$ via the natural map $M\to M\otimes_{A}\widehat{A}$. This map becomes an isomorphism upon tensoring with $\widehat{A}$, hence was an isomorphism). The result thus follows from the proposition.
\end{proof}

\par\endgroup

\section{Local cohomology on a prescheme}\label{sec:5.4}

Let $X$ be a prescheme, $Y$ a closed subset, and $F$ a sheaf of $\sO_X$-modules. We interpret the local cohomology groups $H_Y^i(F)$ in terms of $\operatorname{Ext}$'s.

For each $n\geq 1$, let $Y_n$ be the subscheme of $X$ defined by the sheaf of ideals $\sJ_Y^{\,n}$, where $\sJ_Y$ is the sheaf of ideals of some scheme structure on $Y$. Then for each $n$ we have a natural map
\[
\operatorname{Hom}_{\sO_X}(\sO_{Y_n}, F)\longrightarrow \Gamma_Y(F),
\]
since an element of the first group is given by a global section of $F$ which is annihilated by $\sJ_Y^{\,n}$, and hence has support on $Y$. Taking the direct limit as $n$ varies, we deduce a map
\[
\varinjlim_n \operatorname{Hom}(\sO_{Y_n}, F)\longrightarrow \Gamma_{Y}(F).
\]
Taking derived functors, we deduce a map of functors from $\DD^{+}(X)$ to $\DD^{+}(\Ab)$
\begin{equation}\tag{*}\label{eq:5.4.support-comparison}
\RR\varinjlim_n \operatorname{Hom}(\sO_{Y_n}, F^{\bullet})\longrightarrow \RR\Gamma_{Y}(F^{\bullet}).
\end{equation}

\begin{booktheorem}{4.1}\label{thm:5.4.1}
If $X$ is a noetherian prescheme, and $F^\bullet\in\DD^+_{\mathrm{qc}}(X)$, i.e., the complex $F^\bullet$ has quasi-coherent cohomology (cf.\ \bookref{sec:2.1}{II,~\S1}), then \eqref{eq:5.4.support-comparison} is an isomorphism.
\end{booktheorem}

\par\medskip\noindent\textbf{Remarks.}
\begin{enumerate}
\item Since there are enough flasque $\sO_X$-modules, $\RR\Gamma_Y(F^\bullet)$ is the same as if one had calculated it in the category $\DD^+(\Ab(X))$.
\item The derived category does not have direct limits, so we cannot write
\[
\varinjlim_n\RHom(\sO_{Y_n}, F^\bullet),
\]
as one is tempted to. However, upon descending from the derived category this difficulty disappears, and we have the Corollary below.
\end{enumerate}
\par\medskip

\begin{bookcorollary}{4.2}\label{cor:5.4.2}
Under the hypotheses of the theorem, the map induced on cohomology,
\[
\varinjlim_n \operatorname{Ext}^{i}(\sO_{Y_{n}}, F^{\bullet})\longrightarrow H_{Y}^{i}(F^{\bullet})
\]
is an isomorphism, for all $i$.
\end{bookcorollary}

\begin{proof}
If $X$ is noetherian, and $F$ quasi-coherent, then the map
\[
\varinjlim_n\operatorname{Hom}(\sO_{Y_{n}}, F)\longrightarrow \Gamma_{Y}(F)
\]
is an isomorphism, since every section of $F$ with support on $Y$ will be annihilated by some power of $\sJ_Y$. Therefore, by the lemma on Way-Out Functors \bookref{prop:1.7.1}{I,~7.1}, the morphism of derived functors \eqref{eq:5.4.support-comparison} is an isomorphism for $F^\bullet\in\DD^+_{\mathrm{qc}}(X)$.

For the Corollary, just observe that taking cohomology of complexes commutes with direct limits.
\end{proof}

\begin{bookcorollary}{4.3}\label{cor:5.4.3}
If $X$ is locally noetherian, and $F^\bullet\in\DD^+_{\mathrm{qc}}(X)$, then the analogous maps
\[
\RR\varinjlim_n \sheafHom(\sO_{Y_n}, F^\bullet)\longrightarrow \RR\underline{\Gamma}_{Y}(F^\bullet)
\]
and
\[
\varinjlim_n \sheafExt^{i}(\sO_{Y_{n}}, F^{\bullet})\longrightarrow \mathcal{H}_{Y}^{i}(F^{\bullet})
\]
are isomorphisms.
\end{bookcorollary}

\begin{proof}
By definition of the derived category, it is sufficient to prove the second. This follows from Corollary~\ref{cor:5.4.2}, since $\sheafExt$ and $\mathcal{H}_{Y}$ are sheaves associated to presheaves given by $\operatorname{Ext}$ and $H_{Y}$.
\end{proof}

\section{Dualizing functors on a local noetherian ring}\label{sec:5.5}

In this section we recall without proof the definition and some properties of dualizing functors on a local noetherian ring $A$, with maximal ideal $\mathfrak{m}$. For proofs, see \cite[\nopp \S4]{LC}.

\begin{bookproposition}{5.1}\label{prop:5.5.1}
Let $I$ be an injective hull of the residue field $k$ of $A$, and denote by $T$ the functor $\operatorname{Hom}(\cdot,I)$. Then for every $A$-module $M$ of finite length, the natural map
\[
M\longrightarrow T T(M)
\]
is an isomorphism.
\end{bookproposition}

\begin{definition*}
A contravariant additive functor $T$ from the category $\mathcal{C}_f$ of $A$-modules of finite length into itself is called a \emph{dualizing functor for $A$ at $\mathfrak{m}$} if there exists an injective hull $I$ of $k$ and an isomorphism of functors $T\cong\operatorname{Hom}(\cdot, I)$.
\end{definition*}

\begin{bookproposition}{5.2}\label{prop:5.5.2}
A contravariant additive functor $T$ from $\mathcal{C}_f$ into itself is a dualizing functor if and only if it is exact, and $T(k)\cong k$. In that case one can take $I=\varinjlim_n T(A/\mathfrak{m}^n)$, and then there is a canonical isomorphism $T\cong\operatorname{Hom}(\cdot, I)$.
\end{bookproposition}

\begin{bookproposition}{5.3}\label{prop:5.5.3}
If $I$ is an injective hull of $k$, then the functor $\operatorname{Hom}(\cdot,I)$ gives an anti-isomorphism of the category (dcc) of $A$-modules with descending chain condition (we call them $A$-modules of co-finite type) and the category (acc) of modules of finite type over the completion of $A$.
\end{bookproposition}

\section{Local duality}\label{sec:5.6}

Throughout this section we let $A$ be a noetherian local ring, with maximal ideal $\mathfrak{m}$, and residue field $k$, and let $R^\bullet$ be a dualizing complex on $X=\operatorname{Spec}A$ (cf.\ section~\hyperref[sec:5.2]{2}). Recall Proposition~\ref{prop:5.3.4} which gives a necessary and sufficient condition for a complex $R^\bullet\in\DD^+_{\mathrm{c}}(A)$ to be dualizing.

\begin{definition*}
We say that the dualizing complex $R^\bullet$ is \emph{normalized} if the integer $d$ of Proposition~\ref{prop:5.3.4} is zero.
\end{definition*}

\par\medskip\noindent\textbf{Remark.}
Since the translate of a dualizing complex is again one, we can normalize by translation.
\par\medskip

\begin{bookproposition}{6.1}\label{prop:5.6.1}
If $R^\bullet$ is normalized, then $\RR\Gamma_x(R^\bullet)$, where $x$ is the closed point of $\operatorname{Spec}A$, is isomorphic in the derived category to an injective hull $I$ of $k$.
\end{bookproposition}

\begin{proof}
The cohomology of $\RR\Gamma_x(R^\bullet)$ is
\[
H_{x}^{i}(R^{\bullet})\cong \varinjlim_n\operatorname{Ext}^{i}(A/\mathfrak{m}^{n},R^{\bullet})
\]
by Corollary~\ref{cor:5.4.2} above. On the other hand, we saw from the proof of Proposition~\ref{prop:5.3.4} that $\operatorname{Ext}^{i}(M,R^\bullet)=0$ for $i\neq0$ and all $M$ of finite length. Furthermore, by Proposition~\ref{prop:5.5.2}, the functor $\operatorname{Ext}^{0}(\cdot,R^\bullet)$ is a dualizing functor, and corresponds to the injective hull of $k$ given by
\[
I = \varinjlim_n \operatorname{Ext}^{0}(A/\mathfrak{m}^{n}, R^{\bullet}).
\]
So we see that $H_{x}^{i}(R^\bullet)=0$ for $i\neq0$, and $H_{x}^{0}(R^\bullet)=I$. Therefore $\RR\Gamma_x(R^\bullet)\cong I$ in the derived category \bookref{prop:1.4.3}{I,~4.3}.
\end{proof}

Now let $M$ be an $A$-module. Then there is a natural homomorphism
\[
\Gamma_x(M)\longrightarrow \operatorname{Hom}_A(\operatorname{Hom}_A(M,R^{\bullet}), \Gamma_x(R^{\bullet})).
\]
This gives rise to a morphism of functors on $\DD^+(A)$,
\[
\RR\Gamma_x(M^{\bullet})\longrightarrow \RHom_A(\RHom_A(M^{\bullet}, R^{\bullet}), \RR\Gamma_x(R^{\bullet})),
\]
since we can take $R^\bullet$ to be an injective complex, in which case $\RR\Gamma_x(R^\bullet)=\Gamma_x(R^\bullet)$ is also injective. Indeed, whenever $J$ is an injective $A$-module, and $Y$ a closed subset of $\operatorname{Spec}A$, then $\Gamma_Y(J)$ is injective (cf.\ \bookref{sec:2.7}{II,~\S7} for the structure of injective $A$-modules). We now assume (without loss of generality) that $R^\bullet$ is normalized, and so, by Proposition~\ref{prop:5.6.1}, have a morphism of functors
\[
\Theta\colon\ \RR\Gamma_x(M^{\bullet})\longrightarrow \operatorname{Hom}_A(\RHom_A(M^{\bullet},R^{\bullet}), I).
\]
(Following the usual conventions, we do not write $\RR$ before a functor which is already exact.)

\begin{booktheorem}[Local Duality]{6.2}\label{thm:5.6.2}
Let $A$ be a local noetherian ring, let $R^\bullet$ be a normalized dualizing complex, let $I$ be the corresponding injective hull of $k$, and let $M^\bullet\in\DD^+_{\mathrm{c}}(A)$ (i.e., the cohomology modules of $M^\bullet$ are of finite type). Then $\Theta$ as defined above is an isomorphism.
\end{booktheorem}

\begin{bookcorollary}{6.3}\label{cor:5.6.3}
With $A$, $R^\bullet$, $I$ as above, let $M$ be a module of finite type. Then the homomorphisms
\[
\Theta^{i}\colon H_{x}^{i}(M)\longrightarrow \operatorname{Hom}_A(\operatorname{Ext}^{-i}(M,R^{\bullet}), I)
\]
induced by $\Theta$ are isomorphisms.
\end{bookcorollary}

\begin{proof}
Using the Lemma on Way-out Functors \bookref{prop:1.7.1}{I,~7.1} one sees that the Corollary is equivalent to the theorem. To prove the Corollary, note that $\Theta^{i}$ is an isomorphism for $M=k$, because $H_{x}^{i}(k)=0$ for $i\neq0$, and $H_{x}^{0}(k)=k$; note also that for all $M$ of finite type, both sides are modules with support at $x$. Thus we are reduced to proving the following

\begin{booklemma}{6.4}\label{lem:5.6.4}
Let $\Theta^{i}\colon S^{i}\to T^{i}$ be a morphism of covariant cohomological functors on the category of modules over a noetherian ring $A$. Assume
\begin{itemize}
\item[(i)] for every maximal ideal $\mathfrak{m}\subseteq A$, $\Theta^{i}(A/\mathfrak{m})$ is an isomorphism, and
\item[(ii)] for every non-maximal prime ideal $\mathfrak{p}\subseteq A$, $S^{i}(A/\mathfrak{p})$ and $T^{i}(A/\mathfrak{p})$ have support $<\operatorname{Supp}A/\mathfrak{p}$.
\end{itemize}
Then $\Theta^{i}(M)$ is an isomorphism for every $A$-module $M$ of finite type.
\end{booklemma}

\begin{proof}
By assumption, $\Theta^{i}(A/\mathfrak{m})$ is an isomorphism for every maximal ideal $\mathfrak{m}$. If $M$ is any module of finite type, then $M$ admits a finite filtration each of whose quotients is of the form $A/\mathfrak{p}$ with $\mathfrak{p}$ prime. Thus by the $5$-lemma we reduce to the case $M=A/\mathfrak{p}$ with $\mathfrak{p}$ not maximal. By noetherian induction, we may assume that $\Theta^{i}(M')$ is an isomorphism for every $M'$ of finite type with $\operatorname{support}<\operatorname{Supp}(A/\mathfrak{p})$.

Now for each $f\in M=A/\mathfrak{p}$, $f\neq0$, we consider the exact sequence
\[
0\longrightarrow M\xrightarrow{f} M\longrightarrow M/fM\longrightarrow 0,
\]
and apply the functors $S^{i}$, $T^{i}$ to it, splitting it as follows
\[
\begin{tikzcd}
	% Top Row
	0 \arrow[r] & K_f \arrow[r] \arrow[d, "\alpha"] & S^i(M) \arrow[r, "f"] \arrow[d, "\beta"] & S^i(M) \arrow[r] \arrow[d] & \dots \\
	% Bottom Row
	0 \arrow[r] & L_f \arrow[r] & T^i(M) \arrow[r, "f"] & T^i(M) \arrow[r] & \dots
\end{tikzcd}
\]
and
\[
\begin{tikzcd}
	% Top Row
	\dots \arrow[r] & S^{i-1}(M) \arrow[r] \arrow[d, "\gamma"'] & S^{i-1}(M/fM) \arrow[r] \arrow[d, "\delta"'] & K_f \arrow[r] \arrow[d, "\alpha"'] & 0 \\
	% Bottom Row
	\dots \arrow[r] & T^{i-1}(M) \arrow[r] & T^{i-1}(M/fM) \arrow[r] & L_f \arrow[r] & 0
\end{tikzcd}
\]
Now by our induction hypothesis, $\delta$ is an isomorphism, so $\alpha$ is surjective. Since $\operatorname{Supp}S^{i}(M)<\operatorname{Supp}M$ by our assumption (ii), each element of $S^{i}(M)$ is annihilated by some non-zero $f\in M$. Thus $S^{i}(M)=\bigcup K_f$ as $f$ ranges over the non-zero elements of $M$. Similarly $T^{i}(M)=\bigcup L_f$. Therefore $\beta$ is surjective. This is true for each $i$, so also $\gamma$ is surjective. But that implies $\alpha$ is injective, so $\alpha$ is an isomorphism. Then, as above, $\beta$ is also an isomorphism.
\end{proof}
\renewcommand{\qedsymbol}{}
\end{proof}

\begin{bookcorollary}{6.5}\label{cor:5.6.5}
For $M$ of finite type, the modules $H_x^{i}(M)$ are of co-finite type (see Proposition~\ref{prop:5.5.3}), and we have a functorial isomorphism
\[
\operatorname{Ext}^{i}(M,R^{\bullet})^{\wedge}\longrightarrow \operatorname{Hom}_A(H_{x}^{-i}(M), I)
\]
of $\widehat{A}$-modules.
\end{bookcorollary}

\begin{proof}
Apply the functor $\operatorname{Hom}_A(\cdot,I)$ to the isomorphism $\Theta^{i}$, and use the fact that for any $A$-module of finite type $M$, this functor applied twice gives the completion $\widehat{M}$ \cite[61]{LC}.
\end{proof}

\par\medskip\noindent\textbf{Remark.}
If $A$ is a regular local ring, then $A$ itself is a dualizing complex for $A$ (see Example~\ref{ex:5.2.2}), and we recover the old local duality theorem \cite[\nopp 6.3]{LC}. This is also true more generally for Gorenstein local rings (Theorem~\hyperref[thm:5.9.1]{9.1} below).
\par\medskip

\section{Application to dualizing complexes}\label{sec:5.7}

Let $X$ be a locally noetherian prescheme, and let $R^\bullet$ be a dualizing complex on $X$. Then for each point $x\in X$, $R_x^\bullet$ is a dualizing complex for the local ring $\sO_{X,x}$ (Corollary~\ref{cor:5.2.3}). Thus by Proposition~\ref{prop:5.3.4} we can find an integer $d(x)$ such that
\[
\operatorname{Ext}_{\sO_{X,x}}^{i}(k(x), R_{x}^{\bullet}) = \begin{cases} 0 & \text{for } i\neq d(x) \\ k(x) & \text{for } i = d(x). \end{cases}
\]

\begin{bookproposition}{7.1}\label{prop:5.7.1}
With the hypotheses above, if $x\rightsquigarrow y$ is an immediate specialization, then
\[
d(y)=d(x)+1.
\]
\end{bookproposition}

\begin{proof}
Since the question is local around $y$, we may assume that $X=\operatorname{Spec}\sO_{X,y}$. Let $Z$ be the reduced induced subscheme structure on $\overline{\{x\}}$, and let $j\colon Z\to X$ be the inclusion. Then $j^{\flat}R^{\bullet}$ is dualizing on $Z$ by Proposition~\ref{prop:5.2.4}, and for any $\sO_Z$-module $F$ we have
\[
\operatorname{Ext}_{\sO_Z}^{i}(F, j^{\flat}R^{\bullet}) = \operatorname{Ext}_{\sO_X}^{i}(j_{*}F, R^{\bullet})
\]
by duality for $j$ \bookref{thm:3.6.7}{III,~6.7}. Thus $d(x),d(y)$ are the same if calculated on $Z$ with $j^{\flat}R^{\bullet}$, and so we reduce to the case $X=Z$ integral local of dimension $1$, with generic point $x$ and closed point $y$.

By translating, we may assume that $R^\bullet$ is normalized, i.e., $d(y)=0$. To calculate $d(x)$ we consider $\operatorname{Ext}_{k(x)}^{i}(k(x),R_x^{\bullet})=\operatorname{Ext}_{A}^{i}(A,R^{\bullet})\otimes_{A} k(x)$, where $X=\operatorname{Spec}A$. Now by the local duality theorem (Theorem~\ref{thm:5.6.2}), $\operatorname{Ext}_{A}^{i}(A,R^{\bullet})$ is dual to $H_{y}^{-i}(A)$. Since $A$ is a domain, $H_{y}^{0}(A)=0$. For $p\neq0,1$, $H_{y}^{p}(A)=0$ since $A$ is of dimension $1$ \cite[\nopp 1.12]{LC}. Thus $\operatorname{Ext}^{i}(A,R^{\bullet})=0$ for $i\neq-1$, and so $d(x)=-1$, as required.
\end{proof}

\begin{definition*}
An integer-valued function $d$ on a prescheme $X$ with the property of the proposition will be called a \emph{codimension function} on $X$.
\end{definition*}

\begin{definition*}
A prescheme $X$ is \emph{catenary} if whenever $Z\subseteq Z'$ are irreducible closed subsets of $X$, then every maximal chain
\[
Z=Z_0<Z_1<\cdots<Z_n=Z'
\]
of irreducible closed subsets between $Z$ and $Z'$ has the same length $n$.
\end{definition*}

\begin{bookcorollary}{7.2}\label{cor:5.7.2}
A locally noetherian prescheme $X$ which admits a dualizing complex is catenary, and has finite Krull dimension.
\end{bookcorollary}

\begin{proof}
Indeed, the existence of a codimension function $d$ on $X$ implies that $X$ is catenary. On the other hand, since $R^\bullet$ is quasi-isomorphic to a bounded complex of injective sheaves, the function $d$ is bounded in both directions, so $X$ has finite Krull dimension.
\end{proof}

\begin{definition*}
Let $X$ be a locally noetherian prescheme, and let $R^\bullet$ be a dualizing complex on $X$. We define the \emph{filtration} $Z^\bullet$ associated to $R^\bullet$ by
\[
Z^{p} = \{x\in X\mid d(x)\geq p\}
\]
for each $p$. Then by virtue of the previous proposition and corollary, $Z^\bullet$ is a finite filtration of $X$, and each $x\in Z^{p}-Z^{p+1}$ is maximal.
\end{definition*}

\begin{bookproposition}{7.3}\label{prop:5.7.3}
Let $X$ be a locally noetherian prescheme, let $R^\bullet$ be a dualizing complex on $X$, and let $Z^\bullet$ be the associated filtration. Then $R^\bullet$ is Gorenstein with respect to $Z^\bullet$ (cf.\ \bookref{sec:4.3}{IV,~\S3}). Moreover, for each $x\in X$, the injective hull $J(x)$ of $k(x)$ over $\sO_{X,x}$ occurs exactly once as a direct summand of the (injective) complex $E(R^\bullet)$.
\end{bookproposition}

\begin{proof}
Using condition (ii) of \bookref{prop:4.3.1}{IV,~3.1} and \bookref{motif:4.1.F}{IV,~1.F}, it is sufficient to show that
\[
H_{x}^{i}(R^{\bullet}) = \begin{cases} 0 & \text{for } i\neq d(x) \\ J(x) & \text{for } i = d(x). \end{cases}
\]
But this follows from Proposition~\ref{prop:5.6.1} above.
\end{proof}

\section{Pointwise dualizing complexes and $f^{\sharp}$}\label{sec:5.8}
\begingroup
\setlength{\emergencystretch}{2em}

In this section we give a generalization of the notion of dualizing complex, which may exist on a locally noetherian prescheme of infinite Krull dimension. We apply this notion to show that $f^{\sharp}$ of a dualizing complex is dualizing, where $f$ is a smooth morphism. Throughout this section, $X$ will denote a locally noetherian prescheme.

\begin{definition*}
A complex $R^\bullet\in\DD_{\mathrm{c}}^+(X)$ has \emph{pointwise finite injective dimension} (pfid) if for every $x\in X$, $R_x^\bullet\in\DD_{\mathrm{c}}^+(\Spec\sO_{X,x})$ has finite injective dimension.
\end{definition*}

\begin{bookproposition}{8.1}\label{prop:5.8.1}
Let $R^\bullet\in\DD_{\mathrm{c}}^+(X)$ have pointwise finite injective dimension. Then conditions (i)--(iv) of Proposition~\ref{prop:5.2.1} are equivalent.
\end{bookproposition}

\begin{proof}
The same as loc.\ cit.\ since one can check reflexivity pointwise.
\end{proof}

\begin{definition*}
A complex $R^\bullet\in\DD^+_{\mathrm{c}}(X)$ with pfid is \emph{pointwise dualizing} if it satisfies the equivalent conditions of the proposition.
\end{definition*}

\par\medskip\noindent\textbf{Remarks.}
\begin{enumerate}
\item Since fid implies pfid, we see that any dualizing complex is pointwise dualizing.
\item If $X=\Spec A$ where $A$ is a local ring, then $\text{fid}\iff\text{pfid}$, so a complex is dualizing if and only if it is pointwise dualizing.
\item We say a complex $F^\bullet\in\DD(X)$ is \emph{pointwise bounded below} if for every $x\in X$, $F_x^\bullet\in\DD^+(\Spec\sO_{X,x})$, and we write $F^\bullet\in\DD^{\mathrm{p}+}(X)$. Similarly we define $\DD^{\mathrm{p}-}(X)$ and $\DD^{\mathrm{pb}}(X)$. These are full subcategories of $\DD(X)$ (cf.\ \bookref{sec:1.4}{I,~\S4}, \bookref{sec:1.5}{I,~\S5}). Now if $R^\bullet\in\DD_{\mathrm{c}}^+(X)$ is pointwise dualizing, then the functor $\RsheafHom^\bullet(\cdot,R^\bullet)$ gives an anti-isomorphism of $\DD_{\mathrm{c}}^{\mathrm{p}-}(X)$ and $\DD_{\mathrm{c}}^{\mathrm{p}+}(X)$, and an auto-anti-isomorphism of $\DD_{\mathrm{c}}^{\mathrm{pb}}(X)$.
\item If $R^\bullet$ is a pointwise dualizing complex on $X$, we can define as in section~\hyperref[sec:5.7]{7} the associated codimension function $d\colon X\to\mathbb{Z}$, and the associated filtration $Z^\bullet$. Proposition~\ref{prop:5.7.1} holds in this case, so the filtration $Z^\bullet$ satisfies the conditions of \bookref{sec:4.3}{IV,~\S3}, in particular, it is separated. Thus we see as in Corollary~\ref{cor:5.7.2} that $X$ is catenary.
\item If $R^\bullet$ is pointwise dualizing on $X$, and if $R^\bullet\in\DD_{\mathrm{c}}^{\mathrm{b}}(X)$, then $R^\bullet$ is Gorenstein with respect to the filtration $Z^\bullet$ (cf.\ Proposition~\ref{prop:5.7.3}). The boundedness hypothesis is necessary for \bookref{prop:4.3.1}{IV,~3.1}.
\end{enumerate}
\par\medskip

\begin{example*}
If $X$ is a regular prescheme, then $\sO_X$ is a pointwise dualizing complex for $X$. Indeed, for each $x\in X$, $\Spec\sO_{X,x}$ is a regular prescheme of finite Krull dimension, so we apply Example~\ref{ex:5.2.2} above.
\end{example*}

\begin{bookproposition}{8.2}\label{prop:5.8.2}
Let $R^\bullet\in\DD_{\mathrm{c}}^{+}(X)$ be pointwise dualizing, and assume that $X$ is noetherian and of finite Krull dimension. Then $R^\bullet$ has finite injective dimension and so is dualizing (by Corollary~\ref{cor:5.2.3}).
\end{bookproposition}

\begin{proof}
Since $X$ is noetherian and of finite Krull dimension, the codimension function $d$ associated to $R^\bullet$ is bounded. But we saw in the proof of Proposition~\ref{prop:5.3.4} that for every $x\in X$, and every module $M$ of finite type over $\sO_{X,x}$,
\[
\operatorname{Ext}_{\sO_{X,x}}^{i}(M, R_{x}^{\bullet}) = 0 \qquad\text{for } i\notin[d(x)-\dim M, d(x)].
\]
Thus in particular for every coherent sheaf $F$ on $X$,
\[
\sheafExt_{X}^{i}(F,R^{\bullet}) = 0 \qquad\text{for } i > \sup_{x\in X} d(x),
\]
and so $R^\bullet$ has finite injective dimension \bookref{prop:2.7.20}{II,~7.20 $(\underline{\mathrm{iii}})_{\mathrm{c}}$}.
\end{proof}

\begin{booktheorem}{8.3}\label{thm:5.8.3}
Let $f\colon X\to Y$ be a smooth morphism of locally noetherian preschemes whose fibres have bounded dimension,\footnote{This hypothesis is added in the author\textquotesingle s errata to ensure that $f^\sharp(R^\bullet)$ is bounded below.} and let $R^{\bullet}\in\DD^{+}_{\mathrm{c}}(Y)$ be pointwise dualizing. Then $f^{\sharp}(R^{\bullet})$ (cf.\ \bookref{sec:3.2}{III,~\S2}) is pointwise dualizing on $X$. If furthermore $Y$ is noetherian, $f$ of finite type, and $R^{\bullet}$ dualizing, then $f^{\sharp}(R^{\bullet})$ is also dualizing.
\end{booktheorem}

\begin{proof}
For the first statement, the question is local on $X$ and $Y$, so we reduce immediately to the case $Y=\Spec A$, with $A$ a local ring, and $X$ of finite type, affine over $Y$. Then $X$ admits a closed immersion into a suitable affine space $\mathbf{A}_{Y}^{n}$, so using \bookref{prop:3.8.4}{III,~8.4} and Proposition~\ref{prop:5.2.4} (which is valid also for ``pointwise dualizing''), we reduce to the case $X=\mathbf{A}_{Y}^{n}$.

\medskip\noindent\textit{Case 1.} $A$ is a regular local ring. Then $R^\bullet$ is dualizing, since $A$ is local, and by uniqueness (Theorem~\ref{thm:5.3.1}) we may assume, after shifting if necessary, that $R^\bullet\cong A$. Then $f^{\sharp}(R^\bullet)\cong\omega_{X/Y}[n]=\Omega_{X/Y}^{n}[n]$, which is a shift of a locally free sheaf of rank one, and hence pointwise dualizing (Example~\ref{ex:5.2.2}, since $X$ is regular if $Y$ is).

\medskip\noindent\textit{Case 2.} $A$ is a quotient of a regular local ring $A'$. Let $Y'=\Spec A'$. Then $X$ is obtained from the morphism $\mathbf{A}_{Y'}^{n}\to Y'$ by the base extension $Y\to Y'$, and we reduce to the previous case by means of \bookref{cor:3.6.4}{III,~6.4}, Proposition~\ref{prop:5.2.4}, and the uniqueness (Theorem~\ref{thm:5.3.1}) on $Y$.

\medskip\noindent\textit{General case.} $A$ is an arbitrary noetherian local ring.

Let $\widehat{A}$ be the completion of $A$, and let $Y'=\Spec\widehat{A}$. We make the base extension $u\colon Y'\to Y$, and then $u^{*}(R^\bullet)$ is dualizing on $Y'$ by Corollary~\ref{cor:5.3.5}. But $\widehat{A}$ is a quotient of a regular local ring by Cohen's structure theorem \cite[\nopp (31.1)]{14}, so $f'^{\sharp}(u^{*}(R^\bullet))$ is dualizing on $X' = X\times_{Y} Y'$, by Case 2. Now for any point $x\in X$, the completion $\widehat{\sO}_{X,x}$ of its local ring is also the completion of the ring $\sO_{X,x}\otimes_{A}\widehat{A}$, which is the semilocal ring of the points of $X'$ lying over $x$, and so applying Corollary~\ref{cor:5.3.5} once in each direction, and \bookref{prop:3.2.1}{III,~2.1} we see that $f^{\sharp}(R^\bullet)$ is pointwise dualizing on $X$.

Now if $Y$ is noetherian, $f$ of finite type, and $R^\bullet$ dualizing, then $Y$ has finite Krull dimension by Corollary~\ref{cor:5.7.2}, hence $X$ has finite Krull dimension, and so $f^{\sharp}(R^\bullet)$ is dualizing by Proposition~\ref{prop:5.8.2}.
\end{proof}

\begin{bookcorollary}{8.4}\label{cor:5.8.4}
Under the hypotheses of the theorem, let $d_{Y}$ and $d_{X}$ be the codimension functions associated with the pointwise dualizing complexes $R^{\bullet}$ and $f^{\sharp}(R^{\bullet})$, respectively. Then for each $x\in X$,
\[
d_{X}(x) = d_{Y}(y) - \operatorname{tr.d.} k(x)/k(y)
\]
where $y=f(x)$.
\end{bookcorollary}

\begin{proof}
Chasing through the reductions above (and noting that this formula holds trivially for a finite morphism and $f^{\flat}(R^{\bullet})$), we reduce to the case $Y=\Spec A$, $A$ a regular local ring, $R^\bullet=A$, and $X=\mathbf{A}_{Y}^{n}$. Then $f^{\sharp}(R^\bullet)=\omega_{X/Y}[n]$, a shift of a locally free sheaf of rank one, so for any $x\in X$ we have $d_X(x)=\dim\sO_{X,x}-n$, and for any $y\in Y$, $d_Y(y)=\dim\sO_{Y,y}$ \cite[\nopp \S4]{LC}. Our formula is then the usual dimension formula for a local homomorphism of local rings (cf.\ \cite[\nopp $0_{\mathrm{IV}}$,~17.3.3]{EGA}).
\end{proof}

\par\medskip\noindent\textbf{Remark.}
Using Proposition~\ref{prop:5.2.4}, the result of the theorem can be extended to show that for any embeddable morphism $f$ of noetherian preschemes \bookref{sec:3.8}{III,~\S8}, $f^{!}$ takes dualizing complexes to dualizing complexes. The formula of the Corollary also extends to this case.
\par\medskip

\begin{bookproposition}{8.5}\label{prop:5.8.5}
Let $f\colon X\to Y$ be an embeddable morphism of locally noetherian preschemes, let $R^\bullet\in\DD^+_{\mathrm{c}}(Y)$ be a dualizing complex on $Y$, and let $\mathscr{D}$ denote the functor $\RsheafHom^\bullet_{Y}(\cdot, R^\bullet)$ or $\RsheafHom^\bullet_{X}(\cdot, f^{!}R^\bullet)$ (the one meant will be clear from the context). Then there is a functorial isomorphism
\[
f^{!}(F^{\bullet}) \cong \mathscr{D}(\LL f^{*}(\mathscr{D}(F^{\bullet})))
\]
for all $F^\bullet\in\DD_{\mathrm{c}}^{+}(Y)$.
\end{bookproposition}

\begin{proof}
Apply \bookref{prop:3.8.8}{III,~8.8(7)} to $\mathscr{D}(F^\bullet)$ and $R^\bullet$.
\end{proof}

\begin{bookcorollary}{8.6}\label{cor:5.8.6}
Under the hypotheses of the proposition, assume furthermore that $f$ is flat, of finite type, and that $Y$ is noetherian. Then there is a functorial isomorphism
\[
f^{!}(F^{\bullet}\derivedtensor G^{\bullet}) \xrightarrow{\sim} f^{!}(F^{\bullet}) \derivedtensor f^{*}(G^{\bullet})
\]
for $F^\bullet\in\DD_{\mathrm{c}}^{+}(Y)$ and $G^\bullet\in\DD_{\mathrm{c}}^{\mathrm{b}}(Y)_{\mathrm{fTd}}$.
\end{bookcorollary}

\begin{proof}
Note first that Theorem~\ref{thm:5.8.3} and the Remark following apply to show that $f^{!}(R^\bullet)$ is dualizing on $X$. The result now follows from the Proposition, using \bookref{prop:2.5.8}{II,~5.8} and Proposition~\ref{prop:5.2.6} above.
\end{proof}

\par\endgroup

\section{Gorenstein preschemes}\label{sec:5.9}
\begingroup
\setlength{\emergencystretch}{2em}

\begin{booktheorem}{9.1}\label{thm:5.9.1}
Let $A$ be a noetherian local ring. Then the following conditions are equivalent:
\begin{itemize}
\item[(i)] $A$ is a dualizing complex for itself.
\item[(ii)] $A$ has a finite injective resolution.
\item[(iii)] $A$ is Cohen-Macaulay, and whenever $x_1,\dots,x_n$ is a maximal $A$-sequence, then the ideal $(x_1,\dots,x_n)$ is irreducible.
\item[(iv)] There is an integer $d$ such that
\[
\operatorname{Ext}^{i}(k,A) = \begin{cases} 0 & \text{for } i\neq d \\ k & \text{for } i = d \end{cases}
\]
where $k$ is the residue field of $A$.
\item[(v)] $A$ is Gorenstein with respect to a suitable filtration $Z^{\bullet}$ of $\Spec A$ (cf.\ \bookref{sec:4.3}{IV,~\S3}).
\item[(vi)] There is an integer $d$ such that
\[
H_{x}^{i}(A) = \begin{cases} 0 & \text{for } i\neq d \\ \text{is injective} & \text{for } i = d \end{cases}
\]
where $x$ is the closed point of $\Spec A$.
\end{itemize}
\end{booktheorem}

\begin{proof}
(i) $\Longrightarrow$ (ii) By definition of dualizing complex.

\noindent (ii) $\Longleftrightarrow$ (iii) $\Longleftrightarrow$ (iv) Proved by Bass \cite[\nopp Thm.~4.1]{2}. This article, incidentally, is a good summary of all previous occurrences of Gorenstein rings in the literature.

\noindent (iv) $\Longrightarrow$ (i) This is Proposition~\ref{prop:5.3.4} above.

\noindent (i) $\Longrightarrow$ (v) This is Proposition~\ref{prop:5.7.3} above.

\noindent (v) $\Longrightarrow$ (vi) By \bookref{prop:4.3.4}{IV,~3.4}, $E(A)$ is an injective resolution of $A$, and furthermore, for each $p\in\mathbb{Z}$, $E^{p}(A)$ is a direct sum of constant (injective) sheaves on subspaces $\overline{\{y\}}$ of $\Spec A$ with $y\in Z^{p}-Z^{p+1}$. We can calculate $H_{x}^{i}(A)$ as $H^{i}(\Gamma_{x}(E(A)))$. But $\Gamma_{x}(E(A))$ by what we have just seen, consists of a single injective sheaf in some degree (say $d$), since $x$ is a minimal point of $\Spec A$.

\noindent (vi) $\Longrightarrow$ (iv) First note by \cite[\nopp 6.4, part~4]{LC} that $H_{x}^{n}(A)\neq0$, where $n=\dim A$, so $n=d$. (The hypothesis in loc.\ cit.\ that $A$ is a quotient of a regular local ring can be circumvented by passing to the completion \cite[\nopp 5.9]{LC}.) Therefore $A$ is Cohen-Macaulay \cite[\nopp 3.10]{LC}. Let $(x_1,\dots,x_n)$ be a maximal $A$-sequence, and let $N=(x_1,\dots,x_n)$. Then $H_{x}^{n}(A)$ is an essential extension of $A/N$ \cite[\nopp proof of Prop.~2]{11}, and so is irreducible, since $A/N$ is a simple $A$-module, and hence irreducible. We conclude that $H_{x}^{n}(A)$ is an irreducible injective $A$-module with support at $x$, and so is an injective hull of $k$ \bookref{prop:2.7.4}{II,~7.4}.

Now for any $A$-module $M$,
\[
\operatorname{Hom}(k,M) = \operatorname{Hom}(k, \Gamma_{x}(M)).
\]
Furthermore, the functor $\Gamma_{x}$ takes injectives into injectives, so we have a spectral sequence of derived functors, which degenerates in this case to give
\[
\operatorname{Ext}^{i}(k,A) \cong \operatorname{Hom}(k, H_{x}^{i}(A)) = \begin{cases} 0 & i\neq n \\ k & i = n. \end{cases}
\]
\end{proof}

\begin{definition*}
A local noetherian ring satisfying the equivalent conditions of the theorem is called a (local) \emph{Gorenstein ring}.
\end{definition*}

\par\medskip\noindent\textbf{Remark.}
For a local ring, regular $\Longrightarrow$ complete intersection $\Longrightarrow$ Gorenstein $\Longrightarrow$ Cohen-Macaulay, and these implications are all strict \cite{2}.
\par\medskip

\begin{definition*}
A prescheme is \emph{Gorenstein} if all of its local rings are Gorenstein local rings.
\end{definition*}

\begin{bookcorollary}{9.2}\label{cor:5.9.2}
A localization of a Gorenstein local ring is Gorenstein.
\end{bookcorollary}

\begin{proof}
Follows from condition (i) of the Theorem, and Corollary~\ref{cor:5.2.3} above.
\end{proof}

\begin{bookproposition}{9.3}\label{prop:5.9.3}
Let $f\colon X\to\Spec k$ be an embeddable morphism, where $k$ is a field. Then $X$ is Gorenstein if and only if $f^{!}(k)$ is locally isomorphic in $\DD^+(X)$ to $L[n]$, where $L$ is an invertible sheaf and $n$ is an integer.
\end{bookproposition}

\begin{proof}
Indeed, by the remark following Corollary~\ref{cor:5.8.4}, $f^{!}(k)$ is a dualizing complex on $X$. But $X$ is Gorenstein if and only if $\sO_X$ is a dualizing complex on $X$. So the result follows from the uniqueness of the dualizing complex (Theorem~\ref{thm:5.3.1} above).
\end{proof}

\begin{bookcorollary}{9.4}\label{cor:5.9.4}
Let $X$ be a Gorenstein prescheme of finite type over a field $k$, and let $k\subseteq K$ be a field extension. Then $X\otimes_k K$ is also Gorenstein.
\end{bookcorollary}

\begin{proof}
The question is local on $X$, and $f^{!}$ is compatible with flat base extension \bookref{thm:3.8.7}{III,~8.7(5)}.
\end{proof}

\begin{bookcorollary}{9.5}\label{cor:5.9.5}
Let $K,L$ be extension fields of a field $k$, and assume that one of them is finitely generated. Then $K\otimes_k L$ is a Gorenstein ring (i.e., every localization of it is a local Gorenstein ring).
\end{bookcorollary}

\begin{proof}
Similar to the above. One can also give a direct proof, by induction on the number of generators. This strengthens \cite[\nopp IV,~6.7.1.1]{EGA} which says that $K\otimes_k L$ is Cohen-Macaulay.
\end{proof}

\begin{bookproposition}{9.6}\label{prop:5.9.6}
Let $f\colon X\to Y$ be a flat morphism of locally noetherian preschemes, with $Y$ Gorenstein (but no finiteness assumption on $f$). Then $X$ is Gorenstein if and only if all the fibres $X_{y}$, for $y\in Y$, are Gorenstein preschemes.
\end{bookproposition}

\begin{proof}
The question is local, so we reduce to the following: let $A\to B$ be a local homomorphism of local rings, with $A$ Gorenstein, and $B$ flat over $A$. Then $B$ is Gorenstein if and only if $B/\mathfrak{m}_A B$ is Gorenstein.

To prove this, we use condition (iv) of Theorem~\ref{thm:5.9.1}. By flatness,
\[
\operatorname{Ext}_{B}^{i}(B/\mathfrak{m}_A B, B) = \begin{cases} 0 & \text{for } i\neq d \\ B/\mathfrak{m}_A B & \text{for } i = d \end{cases}
\]
for a suitable $d$. Thus the spectral sequence of $\operatorname{Ext}$'s for the change of ring from $B$ to $B/\mathfrak{m}_A B$ degenerates, and we have
\[
\operatorname{Ext}_{B}^{i}(k_B, B) \cong \operatorname{Ext}_{B/\mathfrak{m}_A B}^{i-d}(k_B, B/\mathfrak{m}_A B)
\]
for each $i$, whence the result.
\end{proof}

\begin{bookexercise}{9.7}\label{rem:5.9.7}
Let $f\colon X\to Y$ be a flat morphism of finite type of locally noetherian preschemes. Show that $f^{!}(\sO_Y)$ (which is defined locally on $X$) is locally isomorphic in $\DD^+(X)$ to an invertible sheaf shifted by an integer if and only if all the fibres $X_y$ of $f$, for $y\in Y$, are Gorenstein preschemes. Also show that, locally on $X$, $f^{!}(\sO_Y)$ has a unique non-zero cohomology sheaf, which is flat over $Y$, if and only if all the fibres of $f$ are Cohen-Macaulay preschemes. Such morphisms are called \emph{Gorenstein morphisms} (resp.\ \emph{Cohen-Macaulay morphisms}).

Compare \cite[\nopp IV,~6.3 and 6.7]{EGA} for statements analogous to Corollary~\ref{cor:5.9.5} and Proposition~\ref{prop:5.9.6} for Cohen-Macaulay.
\end{bookexercise}

\par\endgroup

\section{Existence of Dualizing Complexes}\label{sec:5.10}

We now draw together the experience gained in this chapter to make some remarks about the existence of a dualizing complex on a locally noetherian prescheme $X$.

(We leave the analogous discussion for pointwise dualizing complexes to the reader.)

\medskip\noindent\textbf{Sufficient conditions.}
\begin{enumerate}
\item If $X$ is Gorenstein and of finite Krull dimension, then $\sO_X$ is a dualizing complex for $X$. In particular, any regular prescheme of finite Krull dimension has a dualizing complex.
\item If $f\colon X\to Y$ is a morphism of finite type, with $Y$ noetherian, and if $Y$ admits a dualizing complex $R^\bullet$, we will see in the next chapter that $X$ admits a dualizing complex $f^{!}(R^\bullet)$. (If $f$ is embeddable this result follows from section~\hyperref[sec:5.8]{8} above.) Thus any noetherian prescheme of finite type over a Gorenstein prescheme of finite Krull dimension admits a dualizing complex.
\item In particular, any prescheme of finite type over a field $k$ admits a dualizing complex.
\item If $X$ is the spectrum of a complete local ring $A$, then $X$ admits a dualizing complex. Indeed, $A$ is a quotient of a regular local ring by the Cohen structure theorem \cite[\nopp (31.1)]{14}.
\end{enumerate}

\medskip\noindent\textbf{Necessary conditions.}
\begin{enumerate}
\item If $X$ admits a dualizing complex, then $X$ has finite Krull dimension (Corollary~\ref{cor:5.7.2}).
\item If $X$ admits a dualizing complex, then $X$ is catenary (in fact universally catenary since any prescheme of finite type over $X$ admits a dualizing complex, at least locally).
\item More precisely, $X$ admits a codimension function $d$.
\item If $A$ is a local noetherian ring admitting a dualizing complex, and if $\mathfrak{p}\subseteq A$ is a prime ideal, then $\widehat{A}\otimes_A k(\mathfrak{p})$ is a Gorenstein ring. This follows from the following
\end{enumerate}

\begin{bookproposition}{10.1}\label{prop:5.10.1}
If $A$ is a local noetherian domain, with quotient field $K$, which admits a dualizing complex, then $\widehat{A}\otimes_A K$ is a Gorenstein ring.
\end{bookproposition}

\begin{proof}
Let $R^\bullet$ be a dualizing complex on $A$. Then $R^\bullet\otimes_A\widehat{A}$ is dualizing on $\widehat{A}$ by Corollary~\ref{cor:5.3.5}, and so by localization, $(R^\bullet\otimes_A\widehat{A})\otimes_A K$ is dualizing on $\widehat{A}\otimes_A K$. But $R^\bullet\otimes_A K$ is dualizing on $K$, so by Theorem~\ref{thm:5.3.1} we may assume $R^\bullet\otimes_A K\cong K$ after translating if necessary. But $(R^\bullet\otimes_A\widehat{A})\otimes_A K \cong (R^\bullet\otimes_A K)\otimes_K (\widehat{A}\otimes_A K)$ so $\widehat{A}\otimes_A K$ is a dualizing complex for itself, and we are done (Theorem~\ref{thm:5.9.1} (i)).
\end{proof}

\begin{example*}
There are local noetherian domains of dimension $3$ which are not catenary, and local noetherian domains of dimension $2$ which are not universally catenary, and which therefore have no dualizing complexes \cite[\nopp App.~A1, Ex.~2]{14}.
\end{example*}

\par\medskip\noindent\textbf{Problems.}
\begin{enumerate}
\item We do not know if $\widehat{A}\otimes_A K$ is a Gorenstein ring, even if $A$ is a local domain of dimension $1$, hence we do not know whether every local domain of dimension $1$ admits a dualizing complex.
\item If $X$ is noetherian and of finite Krull dimension, and if it admits a dualizing complex locally, and if it admits a codimension function $d$, we do not know whether $X$ admits a dualizing complex globally.
\end{enumerate}
\par\medskip

\noindent\textit{Author's addendum (errata, original p.~423).}
Giraud~\cite{6} provides an obstruction in $H^2(X,\mathbb{G}_m)$, for the Zariski topology, whose vanishing is necessary and sufficient for a locally existing dualizing complex to exist globally. The author also notes that a noetherian scheme of finite Krull dimension can admit dualizing complexes locally but none globally because it has no global codimension function.
